\graphicspath{{chapters/contents/sessio_02/images/}}

\section{Sessió 02 (21/09/26). Solucions.}


\subsection{El grup de control}

Primer de tot, hem de mesurar la temperatura de les pedres que hem deixat dins de la caixa. Ja sabem que no haurien d’haver canviat gaire de temperatura, però les hem de mesurar igualment. En ciència, donar una cosa per feta abans de mesurar-la acostuma a ser una manera força eficient de començar malament.

Aquestes pedres formen el nostre grup de control. Necessitem una referència amb la qual comparar les pedres que hem exposat al Sol. Aquesta idea apareixerà una vegada i una altra en els experiments: si volem saber què ha provocat un canvi, necessitem alguna cosa que ens permeti veure què hauria passat sense aquest canvi.

Les temperatures obtingudes en el grup de control van ser:

Grup A: 25,4 °C (negra), 25,3 °C (grisa) i 25,6 °C (blanca).

Grup B: 27,7 °C (negra), 25,2 °C (grisa) i 25,2 °C (blanca).

En general, és el que esperàvem: tres pedres guardades dins d’una caixa i sense exposició directa a la llum solar haurien de mantenir temperatures semblants. Al grup A això es veu molt clarament; al grup B, en canvi, la pedra negra presenta una temperatura una mica més alta que les altres. Això ja ens dona una primera pista important: mesurar correctament la temperatura de la pedra negra no serà tan fàcil com semblava.


\subsection{El grup experimental (pedres exposades al Sol)}

Deixem les pedres al Sol i busquem un lloc protegit del vent, perquè el moviment de l’aire podria refredar-les per convecció. També evitem que toquin directament el terra i les col·loquem sobre uns suports, de manera que no guanyin calor per contacte.

Després de 30 minuts d’exposició, revisem què ha passat. En aquest punt considerem que les pedres han arribat aproximadament a l’equilibri tèrmic: ja s’han escalfat tant com permeten les condicions de l’experiment.

\begin{figure}[H]
    \centering
    \includegraphics[width=1\linewidth]{piedras.png}
    \caption{Hem trobat un lloc ben assolellat i protegit del vent per uns farigols molt simpàtics. Imatge creada per l’autor, CC BY 4.0.}
\end{figure}

A les taules següents es recullen les mesures obtingudes per cadascun dels dos grups.

\begin{table}[h]
\centering
\caption{Resultats del grup A}
\begin{tabular}{lccc}
\hline
 & \textbf{Negre} & \textbf{Gris} & \textbf{Blanc} \\
\hline
Control  & 25,4 & 25,3 & 25,6 \\
Mesura 1 & 42,9 & 40,6 & 37,1 \\
Mesura 2 & 38,6 & 39,8 & 39,3 \\
Mesura 3 & 39,6 & 40,4 & 36,7 \\
\hline
Mitjana  & 40,37 & 40,27 & 37,70 \\
\hline
\end{tabular}
\end{table}

\begin{table}[h]
\centering
\caption{Resultats del grup B}
\begin{tabular}{lccc}
\hline
 & \textbf{Negre} & \textbf{Gris} & \textbf{Blanc} \\
\hline
Control  & 27,7 & 25,2 & 25,2 \\
Mesura 1 & 41,0 & 40,0 & 38,0 \\
Mesura 2 & 39,1 & 38,0 & 36,9 \\
Mesura 3 & 38,8 & 38,0 & 36,7 \\
Mesura 4 & 41,3 & 40,8 & 37,3 \\
\hline
Mitjana  & 40,05 & 39,20 & 37,23 \\
\hline
\end{tabular}
\end{table}

És el resultat que esperàvem?

Les pedres exposades al Sol s’han escalfat clarament respecte del grup de control. També observem una tendència general: com més fosc és el color, més alta és la temperatura assolida.

La diferència entre la pedra blanca i les més fosques es veu força bé. En canvi, entre la pedra grisa i la negra la diferència és petita i no podem dir amb gaire seguretat que la negra s’hagi escalfat més.

Per tant, el resultat va en la direcció que esperàvem, però no tan netament com ens hauria agradat.

El grup B ha donat el resultat que esperàvem: la pedra negra s’ha escalfat més que la grisa. Però com poden estar tranquils els científics que això no ha passat simplement per casualitat?

Una possibilitat seria fer moltes més mesures. Però quantes? Aquí la cosa es complica una mica. Cal tenir prou dades, però també disposem d’eines matemàtiques que ens ajuden a decidir si una diferència com la que hem observat podria aparèixer per atzar o si hi ha prou evidència per pensar que l’efecte és real.

Una d’aquestes eines s’anomena contrast d’hipòtesis. Es treballa en cursos més avançats d’estadística i també apareix en algunes modalitats de batxillerat. De moment ens quedem amb la idea important: no n’hi ha prou que el resultat “surti com volíem”; també hem de comprovar fins a quin punt ens en podem refiar.


Com podríem millorar l’experiment?

No hi ha experiments bons o dolents. I, sobretot, no han de sortir perfectes a la primera. Moltes vegades volem posar a prova una idea i el primer muntatge simplement ens serveix per descobrir tot allò que encara no controlem.

Fer un experiment també consisteix en això: repetir, ajustar, controlar millor les condicions i tornar-hi. L’important és aprendre alguna cosa en cada intent. Quan fem una primera prova per veure si una idea funciona i detectar problemes abans de fer un experiment més ben controlat, parlem d’una prova pilot.

En el nostre cas, sembla que hauríem de controlar millor la distància, l’angle i la direcció des d’on mesurem la temperatura. Podríem, per exemple, utilitzar un suport que mantingués sempre el termòmetre en la mateixa posició.

També hi ha el problema de la superfície de les pedres. Les pintures negra i grisa poden reflectir part de la radiació de l’entorn; aquesta propietat està relacionada amb la reflectància i pot afectar la lectura d’un termòmetre infraroig senzill. Finalment, tampoc sabem si la il·luminació solar s’ha mantingut exactament constant durant totes les mesures.

En resum: la física no ens ha fallat. El que encara necessita una mica més de feina és l’experiment.

Per tant, podem millorar l’experiment de diverses maneres:

Assegurar-nos que mesurem sempre la temperatura des de la mateixa distància, amb el mateix angle i des de la mateixa direcció.

Utilitzar una pintura mate, no brillant, per reduir tant com sigui possible l’efecte dels reflexos sobre la mesura.

Fer més mesures i comprovar que no hi hagi tanta diferència entre unes i altres. Si continuen variant molt, probablement hi ha alguna condició que encara no estem controlant prou bé.

